\subsection{Regular expressions matching}
\label{sec:rematch}

\paragraph{Presentation}

In Section~\ref{sec:pmatch}, we explored the matching of simple patterns made of an exact string and possibly a few single-letter wildcards.
We now consider a more complex case with the matching of regular expressions.
These allow for instance for wildcards of any length, or the repetition of a group of letters.
To perform a scan of a reference text and find matching occurrences of an input regular expression, we use the Thompson's construction algorithm~\cite{Thompson1968} that relies on a nondeterministic finite automaton (NFA).

This automaton has transitions labeled by letters but also $\epsilon$-transitions, meaning a state can have several next possible states, hence its nondeterministic nature.
The traversal consists of maintaining a stack of pointers to states in the automaton, and progressing all these pointers at the same time when reading a reference letter, possibly creating new pointers when facing multiple possible transitions.
We also store with the pointers an index to remember at which position in the reference the traversal of this pointer has started.
If a pointer cannot progress further, it gets deleted.
When a pointer however reaches the final state, we record a match at the reference index stored with the pointer.
For each letter in the reference, we add a pointer to the starting state of the automaton to the stack, to try to accept the regular expression at every possible position in the reference.

We implement this algorithm with the BPL library.
The pre-processing (i.e. building the NFA) is performed by the main processor.
The reference text is split over all processing units, with some overlapping between the fragments.
Wildcards of any length thus get limited by the overlapping length, which we set to 128 letters.
Each task receives a reference fragment, the NFA data, and scans its fragment to apply the algorithm and report the hit positions.

We also write an equivalent implementation that uses the main processor only (denoted Std-Multicore), and another that relies on the UPMEM's SDK (denoted SDK-PiM).
 
\paragraph{Limitations encountered}

None: the code can be fully architecture-agnostic.

\paragraph{Benchmarks}

To compare the performances of a given algorithm depending on the target architecture and implementation, we sample a reference genome and generate datasets of regular expressions of various characteristics.
To do so, we pick exact patterns in the genome and alter them to introduce a total of $W$ meta characters such as \texttt{*}, \texttt{+} or \texttt{?}, the original pattern having to be one possible instance of the resulting regular expression.
For instance, \texttt{CACACTGTTCTGTCCCTCCG} may become \texttt{C(AC)*TGT*CTGTCCCTC+G}.
We consider $W$ equal to 0, 5 and 10.
Also, we denote by $L$ the length of the original exact pattern, either 20 or 50 nucleotides.
In a dataset, we aggregate a total $N$ of 10 regular expressions.
We consider two reference genomes: the human genome GRCh38~\cite{HG38} and the human gut microbiome genome UHGG~\cite{UHGG}.

\paragraph{Results} We present the execution time performances of all datasets in \tablename~\ref{tab:rematch_results}.

We see that BPL-Multicore exhibits similar execution times to Std-Multicore for all the GRCh38 datasets.
The overhead appear however a bit less negligible when matching against the UHGG genome which is four times bigger than GRCh38.
On the other hand, BPL-PiM is consistently slightly slower than SDK-PiM.


\begin{table}
    \centering
    \caption{Results of the regular expressions matching benchmarks}
    \label{tab:rematch_results}
    \begin{tabular}{lllllr}
    \toprule
        Reference & $L$ & $N$ & $M$ & Architecture & Thompson NFA (s) \\

\midrule
\multirow{4}{*}{GRCh38} & \multirow{4}{*}{20} & \multirow{4}{*}{10} & \multirow{4}{*}{0} & Std-Multicore & 16.894 \\
&&&& BPL-Multicore & 17.759 \\
\arrayrulecolor{lightgray}\cmidrule{5-6}\arrayrulecolor{black}
&&&& BPL-PiM & 5.068 \\
&&&& SDK-PiM & 4.202 \\
\midrule
\multirow{4}{*}{GRCh38} & \multirow{4}{*}{50} & \multirow{4}{*}{10} & \multirow{4}{*}{0} & Std-Multicore & 17.282 \\
&&&& BPL-Multicore & 18.042 \\
\arrayrulecolor{lightgray}\cmidrule{5-6}\arrayrulecolor{black}
&&&& BPL-PiM & 5.097 \\
&&&& SDK-PiM & 4.217 \\
\midrule
\multirow{4}{*}{GRCh38} & \multirow{4}{*}{20} & \multirow{4}{*}{10} & \multirow{4}{*}{5} & Std-Multicore & 17.733 \\
&&&& BPL-Multicore & 18.602 \\
\arrayrulecolor{lightgray}\cmidrule{5-6}\arrayrulecolor{black}
&&&& BPL-PiM & 4.890 \\
&&&& SDK-PiM & 3.955 \\
\midrule
\multirow{4}{*}{GRCh38} & \multirow{4}{*}{50} & \multirow{4}{*}{10} & \multirow{4}{*}{5} & Std-Multicore & 20.097 \\
&&&& BPL-Multicore & 20.953 \\
\arrayrulecolor{lightgray}\cmidrule{5-6}\arrayrulecolor{black}
&&&& BPL-PiM & 6.142 \\
&&&& SDK-PiM & 4.644 \\
\midrule
\multirow{4}{*}{GRCh38} & \multirow{4}{*}{50} & \multirow{4}{*}{10} & \multirow{4}{*}{10} & Std-Multicore & 19.662 \\
&&&& BPL-Multicore & 20.572 \\
\arrayrulecolor{lightgray}\cmidrule{5-6}\arrayrulecolor{black}
&&&& BPL-PiM & 6.603 \\
&&&& SDK-PiM & 5.095 \\
\midrule
\multirow{4}{*}{UHGG} & \multirow{4}{*}{50} & \multirow{4}{*}{10} & \multirow{4}{*}{10} & Std-Multicore & 83.887 \\
&&&& BPL-Multicore & 87.293 \\
\arrayrulecolor{lightgray}\cmidrule{5-6}\arrayrulecolor{black}
&&&& BPL-PiM & 29.583 \\
&&&& SDK-PiM & 25.709 \\

    \bottomrule
    \end{tabular}
\end{table}

\paragraph{Source code} The implementations are publicly available at \url{https://gitlab.inria.fr/pim/org.pim.rematch-bpl} (GNU AGPLv3).